\hspace{3mm}

\centerline{\bf \Huge Домашнее задание.}
\centerline{\bf \Huge Чилловый разнобой}

\begin{flushright}
    \scriptsize{Не знаю, что тут писать, просто скажу, \\что Алина своровала у меня уже больше 15000 рублей...}
\end{flushright}

\problem{1!} Три целых числа $U, D$ и $E$ таковы, что если $U$ разделить на $D$ получится $E$, а сумма этих чисел равна 2018. Чему могут быть равны эти числа?

\problem{2!} У Германа 8 друзей, и ежедневно в течение 70 дней он отправляет ютуб шортсы 4 из них так, что набор друзей ни в какой день не повторяется. Сколькими способами он может это сделать?

\problem{3} Мачо Баир ежедневно получает 500 сообщений от девочки с химбио и 420 сообщений от девочки с гуманитарного. Известно, что каждое утро ему писала девочка с гуманитарного ровно 300 сообщений, и только после этого начинала писать девочка с химбио. Докажите, что за неделю бурного общения у Мачо Баира было как минимум 7 ситуаций, когда количество сообщений от обеих девочек было одинаковое.

\problem{4}$\bigl[ \textbf{Усиленная КМО Юниоры 3 задача} \bigr]$. Даны $k$ натуральных чисел $a_1, a_2, \dotsc , a_k$. Известно, что $a_1+a_2+\dotsc+a_k = k(p-1)$, где $p \ -$ простое число. Оказалось, что произведение их факториалов 

\centerline{$a_1!\cdot a_2!\cdot \dotsc \cdot a_k!$}

\noindentявляется $k$-й степенью натурального числа. Докажите, что все данные числа равны.